\documentclass[../dayan_kernel_manual.tex]{subfiles}

\begin{document}

\section{执行与内存模型}\label{sec:kernel-010}

\subsection{每核协作式 Runtime}

每个 hart 拥有独立的 CPU-local 状态、init/scheduler 栈、两个 RPC proc 栈、trap 栈
和初始堆。proc 只会在显式调用 yield、等待事件、等待轮询条件、执行阻塞式 RPC 或退出
时切回 scheduler；功能性 CLINT 中断不会抢占正在运行的 proc。

\begin{dayantable}{thread 状态值}{tab:kernel-010}{X[17,l] X[13,l] X[70,l]}
  状态 & 数值 & 语义 \\
  \code{RUNNING} & 0 & 当前正在该 hart 上执行。 \\
  \code{READY} & 1 & 位于 ready queue，等待 scheduler 选择。 \\
  \code{WAIT_POLL} & 2 & scheduler 每轮调用一个不得阻塞的 predicate。 \\
  \code{WAIT_EVENT} & 3 & 已挂入事件所有者的等待链表，由事件显式唤醒。 \\
  \code{WAIT_START} & 4 & Runtime 构造 thread 时使用的短暂内部状态。 \\
  \code{WAIT_EXIT} & 5 & proc 已结束，不会再次入队；当前版本不自动回收 proc 和栈。 \\
\end{dayantable}

\code{dayan_yield} 是算子唯一公开的主动让出接口。\code{proc_wait_poll}、
\code{proc_wait_event} 和 \code{proc_wake} 是 Runtime 内部接口，不在算子符号白名单中。
张量单元等待、resource 等待和 RPC 会在 DAPI/libc 内部选择正确的等待方式。

\subsection{内存域}

\begin{dayantable}{软件可见内存域}{tab:kernel-011}{X[24,l] X[24,l] X[52,l]}
  内存域 & 默认范围/来源 & 使用规则 \\
  固件私有 RAM & 镜像默认从 \code{0x80000000} 开始 & 普通代码、只读数据、每核状态、栈和私有堆；不得假设跨 hart 缓存一致。 \\
  Host coherent window & \code{0x00400000}--\code{0x007fffff} & 4 MiB 硬件一致性访问窗口；具体 allocation 由 Host 动态分配，当前 proc 的基址由 \code{gp/shared} 给出。 \\
  控制空间 & 相对 \code{CONTROL_BASE} & mailbox、interrupt controller、CLINT 等 MMIO，仅由 Runtime/platform 层访问。 \\
\end{dayantable}

系统固件拒绝普通可变 \code{.data}/\code{.bss}，每核可变状态集中在
\code{struct cpu}。coherent window 不属于固件链接布局，链接脚本也拒绝
\code{.coherent} section。Host 根据每个算子提交的用量分配 shared allocation，并在
创建 proc 时把基址写入该 proc 的初始上下文。

算子可以读取全局寄存器变量：

\begin{dayancode}
#include <dayan_sys.h>

void *base = shared;  /* reads gp; do not assign to shared */
\end{dayancode}

\code{shared == NULL} 表示该 proc 没有 shared allocation。该指针只是 allocation
起点；当前 Device ABI 不携带 allocation 长度，精确边界由 Host 和算子布局共同保证。

\subsection{私有堆}

每个 hart 的私有分配器从链接期初始 region 启动，默认大小为 1 MiB。空闲块通过按
大小排序的 intrusive RBTree 做 best-fit 查找，块头中的边界信息用于合并相邻空闲块。
返回地址至少 16-byte 对齐。

初始 region 永不归还。空间不足时，分配器通过 \code{RPC_HEAP_ACQUIRE} 请求 Host
提供新 region；动态 region 完全空闲后通过 \code{RPC_HEAP_RELEASE} 归还。这里的
动态 region 仍是当前 hart 的普通私有堆，不等同于 \code{shared} 指向的跨核一致性
allocation。Host 返回值小于请求、地址非法或与已有 region 重叠时，Runtime 不会接入
该 region，并立即用 release RPC 回滚。接入前会 invalidate 本核旧 cache line；归还
前 clean、Host 确认后 invalidate，避免非一致性缓存污染后续所有者。

\subsection{资源与同步模型}

Dayan resource 使用统一状态机：

\begin{center}
\code{FREE -> BUSY -> READY -> FREE}
\end{center}

\code{alloc} 记录 producer/consumer 数量。producer 数量为 0 时直接进入
\code{READY}；否则最后一次 \code{produce} 完成转换。最后一次 \code{consume}
使资源回到 \code{FREE}。consumer 数量必须非零。

\begin{dayantable}{两类 resource 的并发边界}{tab:kernel-012}{X[24,l] X[33,l] X[43,l]}
  类型 & 存储位置 & 同步方式 \\
  \code{dayan_resource_t} & 当前 hart 的栈或私有堆 & 仅在本 hart 的 proc 间共享；普通 load/store 和核私有事件链表，不生成 AMO/fence。 \\
  \code{dayan_shared_resource_t} & \code{shared} allocation 内 & 可由多个 hart 访问；状态与计数使用原子访问，临界区使用 RV64A 原子锁，等待方轮询一致性状态。 \\
\end{dayantable}

不能把包含 proc 链表节点的 \code{dayan_resource_t} 放入 coherent window，也不能把
\code{dayan_shared_resource_t} 放入普通私有堆后当作跨核对象使用。

\subsection{张量执行单元异步模型}

LSU、MAT、VEC 和 MACRO 单元可以异步执行。对应 wait DAPI 把当前 \code{tp} 作为
Produce token 写入 BASE TCSR 地址 1、2、3、4，并把 proc 挂入事件队列。scheduler
通过 BASE 全局地址 0 反复 pop 已完成 token，再唤醒对应 proc。

每个 wait 占用一个 Produce/Bar 队列项。软件必须遵守硬件容量：全局同时追踪项不超过
32，每个单元不超过 2；超过容量可能造成硬件队列阻塞甚至死锁。

\end{document}
